\onecolumn
%\externaldocument{main}

% ===================================================================
\section*{Appendix}
\renewcommand{\thesubsection}{\Alph{subsection}}

% -------------------------------------------------------------------
% \subsection{Related Work}

% -------------------------------------------------------------------
\subsection{Evaluation Standards}
\textit{\textbf{1) Fixed Point Distance (FPD).}}
This measures the distance of the estimated solution to the minimum of the objective function, where the minimum objective
value is unnecessary to be zero given problem constraints.
%
\begin{align}
\label{eq:fpd}
\textit{LESS}: \quad &\texttt{F}(\mb{y}) = \norm{ \mb{A} \mb{y} - \mb{b} }\ , \\
\textit{IED}: \quad &\texttt{F}(\mb{y}) = \normx{ \mb{y} - \frac{\mb{A} \mb{y}}{\norm{ \mb{A} \mb{y} }} }\ .
\end{align}

\textit{\textbf{2) Eigen Distance.}}
This is intuitive from Eq.~\eqref{eq:eigen_attribute} to evaluate the eigen solution,
\begin{equation}
\label{eq:edist}
D(\mb{y}) = \norm{ \mb{A} \mb{y} - \lambda \mb{y} }\ .
\end{equation}

\textit{\textbf{3) Mean Relative Error (MRE).}}
Given a reference (generally ground truth) $\mb{y}^r$ and $k^{\text{th}}$ sample, the mean relative error is
\begin{equation}
\label{eq:mre}
\texttt{MRE}(\mb{y},\mb{y}^r) = \frac{1}{K} \sum^K_{k=1} \frac{\texttt{F}(\mb{y}_k) - \texttt{F}(\mb{y}^r_k)}{\texttt{F}(\mb{y}^r_k)} \times 100 \%\ ,
\end{equation}
which can be negative when the estimation $\mb{y}$ outperforms the reference $\mb{y}^r$.

% -------------------------------------------------------------------
\subsection{Projected Gradient Descent and Riemannian Manifold Projection}
\begin{figure}[!ht]
\centering
\includegraphics[width=0.18\textwidth]{figures/rie_proj.png}
\caption{\textit{Descent direction projected onto a Riemannian manifold for the sphere constraints.
$-\nabla \texttt{f}(\mb{x}_t)$ is the PGD direction,
$\mb{G}=\text{Proj}_{\texttt{RM}}(-\nabla \texttt{f}(\mb{x}_t))$ is the descent direction on the Riemannian manifold, and $\mb{x}_t$ is on the Riemannian manifold.}}
\label{fig:rie_project}
\end{figure}

% ----
\subsubsection{Projected Gradient Descent}
For ease of optimization, the solution is initialized with the least squares solution without constraints as
\begin{align}
\mb{u}_0 &= \left( \mb{A}^T\mb{A} \right)\inv \mb{A}^T \mb{b}\ , \label{eq:origin} \\
\mb{u}_{0} &\leftarrow \texttt{Proj}_\texttt{Sph}(\mb{u}_{0})\ .
\end{align}
%
Then, the solution is iteratively updated in the gradient descent direction $\Delta \mb{u}_t$ at $t^{\text{th}}$ iteration by following
%
\begin{align}
\mb{u}_{t+1} &= \mb{u}_t +\eta \Delta \mb{u}_t\ , \\
\mb{u}_{t+1} &\leftarrow \texttt{Proj}_{\texttt{Sph}}(\mb{u}_{t+1})
\end{align}
%
with
%
\begin{align}
\Delta \mb{u}_t = -\nabla \texttt{f}(\mb{u}_t)\ ,
\quad
\eta = \frac{\norm{\mb{A}^T (\mb{A}\mb{u}_t-\mb{b})}^2}{\norm{\mb{A}\mb{A}^T(\mb{A}\mb{u}_t-\mb{b})}^2}\ ,
\end{align}
%
where $\eta$ is derived by setting $\nabla \texttt{f}(\eta) \equiv \nabla \texttt{f}(\mb{u}_{t+1})=0$ and $\texttt{Proj}_{\texttt{Sph}}(\mb{u})=\mb{u}/\norm{\mb{u}}$ is the projection onto the unit sphere.
The solution update terminates when the objective value decrease is no greater than a predefined tolerance, $1.0e^{-7}$ in our case, or when it reaches the maximum number of iterations, 100 in our case.

% ----
\subsubsection{Riemannian Manifold Projection}
Since the solution is constrained on the sphere, a simple projection onto the sphere via gradient descent is inefficient considering the inactive orthogonal descent component.
In Fig.~\ref{fig:rie_project}, considering only the descent component on the tangent plane benefits the solution optimality and optimization process, comparing ``PGD" and ``PGD+RM" in Table~\ref{tb:table_1000}.
%
\begin{equation}
\begin{aligned}
\label{eq:P}
\mb{P} = \frac{-\mb{x}_t}{\norm{\mb{x}_t}} \norm{\nabla \texttt{f}(\mb{x}_t)} \text{cos}\theta
= \frac{-\mb{x}_t}{\norm{\mb{x}_t}} \norm{\nabla \texttt{f}(\mb{x}_t)} \frac{\nabla \texttt{f}(\mb{x}_t) \frac{-\mb{x}_t}{\norm{\mb{x}_t}}}{\norm{\nabla \texttt{f}(\mb{x}_t)} \normx{ \frac{\mb{x}_t}{\norm{\mb{x}_t}}}}
= \frac{\mb{x}_t \mb{x}^T_t}{\norm{ \mb{x}_t }^2} \nabla \texttt{f}(\mb{x}_t)\ ,
\end{aligned}
\end{equation}
%
then
\begin{equation}
\begin{aligned}
\label{eq:G}
\mb{G}
= -\nabla \texttt{f}(\mb{x}_t) + \mb{P}
= -\nabla \texttt{f}(\mb{x}_t) + \frac{\mb{x}_t \mb{x}^T_t}{\norm{ \mb{x}_t }^2} \nabla \texttt{f}(\mb{x}_t)\ .
\end{aligned}
\end{equation}

Hence,
%
\begin{equation}
\begin{aligned}
\label{eq:rie_detail}
\texttt{Proj}_{\texttt{RM}}\left( -\nabla \texttt{f}(\mb{x}_t) \right)
= \left( \mb{I}_n - \frac{\mb{x}_t \mb{x}^T_t}{\norm{ \mb{x}_t }^2} \right) \left( -\nabla \texttt{f}(\mb{x}_t) \right)\ ,
\end{aligned}
\end{equation}
where $\mb{x}_t \in \mathbb{R}^n$, $\nabla \texttt{f}(\mb{x}_t) \in \mathbb{R}^n$, $\mb{I}_n$ is an $n \times n$ identity matrix.

% -------------------------------------------------------------------
\subsection{Gradient Transformation with Flipped Eigenvalues}
For different orderings of multiple eigenvalues, for instance, descending-order (concerning the order of eigenvalues) versus ascending-order eigenvalues, the gradients of their associated eigenvectors need to be flipped according to the flipping of the eigenvalues.
Recall the eigen decomposition problem defined in Eq.~\eqref{eq:ed}.
Reserve all of the eigenvalues of $\mb{A}$, that is $n=m$, a flipping function is defined as $\texttt{g}(\mb{x}): \mathbb{R}^{m \times m} \rightarrow \mathbb{R}^{m \times m}$, that is
\begin{equation}
\label{eq:flip}
\texttt{g}(\mb{x}) =
\begin{bmatrix}
 0 & \cdots &  1 \\
 \vdots & \iddots & \vdots \\
1 & \cdots & 0
\end{bmatrix}
%
\mb{x}
%
\begin{bmatrix}
 0 & \cdots &  1 \\
 \vdots & \iddots & \vdots \\
1 & \cdots & 0
\end{bmatrix} \\
=
\mb{I}^{\ast}_{m} \mb{x} \mb{I}^{\ast}_{m}\ ,
\end{equation}
where $(\mb{I}^{\ast}_n \mb{x})$ flips $\mb{x}$ in the vertical direction and $(\mb{x} \mb{I}^{\ast}_n)$ in the horizontal direction.
For the notation simplicity, we denote $\texttt{g}(\mb{x})=\mb{x}^{\text{flip}}$.
Then,
\begin{align}
    \mb{A} \mb{u}
    &= \mb{\Lambda} \mb{u}
    = \mb{I}^{\text{flip}}_m \mb{\Lambda}^{\text{flip}} \mb{I}^{\text{flip}}_m \mb{u}\ , \\
    \mb{I}^{\text{flip}}_m \mb{A} \mb{u}
    &= \mb{\Lambda}^{\text{flip}} \mb{I}^{\text{flip}}_m \mb{u}\ , \\
    \mb{I}^{\text{flip}}_m \mb{A} \left( \mb{I}^{\text{flip}}_m \mb{u}^{\text{flip}} \mb{I}^{\text{flip}}_m \right)
    &= \mb{\Lambda}^{\text{flip}} \mb{u}^{\text{flip}} \mb{I}^{\text{flip}}_m\ , \\
    \mb{A}^{\text{flip}} \mb{u}^{\text{flip}}
    &= \mb{\Lambda}^{\text{flip}} \mb{u}^{\text{flip}}\ ,
\end{align}
where the input $\mb{A} \in \mathbb{R}^{m \times m}$, eigenvectors $\mb{u} \in \mathbb{R}^{m \times m}$, and eigenvalue matrix $\mb{\Lambda} \in \mathbb{R}^{m \times m}$.
It indicates that to obtain the reverse-order eigenvalues, $\mb{A}$ needs to be flipped, which also leads to the flipping of
$\mb{u}$.
Therefore, to compare the descending-order eigen decomposition with an ascending-order one, both $\mb{A}$ and $\mb{u}$ need to be flipped.

For the gradient of loss $L$ over the flipped eigenvectors $\mb{u}^{\text{flip}}$, that is $\nabla_{\mb{u}^{\text{flip}}} L$, since $\mb{u}^{\text{flip}}=\mb{I}^{\text{flip}}_n \mb{u} \mb{I}^{\text{flip}}_n$,
\begin{equation}
    \nabla_{\mb{u}^{\text{flip}}} L
    = \mb{I}^{\text{flip}}_n \nabla_{\mb{u}^{\text{flip}}} L \mb{I}^{\text{flip}}_n
    = \left( \nabla_{\mb{u}} L \right)^{\text{flip}} \ .
\end{equation}

In sum, for the reverse-order eigenvalues, gradients of their associated eigenvectors need to be flipped in both the horizontal and vertical directions, that is to apply Eq.~\eqref{eq:flip}, when the eigenvalue order (or equivalently the eigenvector order) is flipped.

% -------------------------------------------------------------------
% \clearpage
\subsection{Proof and Demo Code for Proposition~\ref{prop:less_hessian}}
We explore the matrix structure by taking an example with $m=2$ and $n=3$, that is $\mb{A}$, $\mb{b}$, and $\mb{y}$ as
%
\begin{equation}
\mb{A} =
\begin{bmatrix}
a_{00} & a_{01} & a_{02} \\
a_{10} & a_{11} & a_{12}
\end{bmatrix}\ ,
\quad
\mb{b} =
\begin{bmatrix}
b_0 \\
b_1
\end{bmatrix}\ , \\
\quad
\mb{y} =
\begin{bmatrix}
y_0 \\
y_1 \\
y_2
\end{bmatrix}\ .
\end{equation}
%
The Jacobian vector and Hessian matrices of the objective function in Eq.~\eqref{eq:obj} over the solution $\mb{y}$ and the Hessian matrix over entries $\mb{A}$ and $\mb{b}$ are
\begin{align}
\nabla_Y \texttt{f}(\mb{A}, \mb{b}, \mb{y})
&= \mb{A}^T \left( \mb{A} \mb{y} - \mb{b} \right) \nonumber \\
&=
\begin{bmatrix}
(a_{00}^2 + a_{10}^2) y_0 + (a_{00} a_{01} + a_{10} a_{11}) y_1 + (a_{00} a_{02} + a_{10} a_{12}) y_2 \\
(a_{01} a_{00} + a_{11} a_{10}) y_0 + (a_{01}^2 + a_{11}^2) y_1 + (a_{01} a_{02} + a_{11} a_{12}) y_2 \\
(a_{02} a_{00} + a_{12} a_{10}) y_0 + (a_{02} a_{01} + a_{12} a_{11} y_1 + (a_{02}^2 + a_{12}^2) y_2)
\end{bmatrix}
-
\begin{bmatrix}
a_{00} b_0+ a_{10} b_1 \\
a_{01} b_0 + a_{11} b_1 \\
a_{02} b_0 + a_{12} b_1
\end{bmatrix}\ ,\\
% ----
\nabla^2_{bY} \texttt{f}(\mb{A}, \mb{b}, \mb{y})
&=
-\begin{bmatrix}
a_{00} & a_{10} \\
a_{01} & a_{11} \\
a_{02} & a_{12}
\end{bmatrix}
=-\mb{A}^T\ , \\
% ----
\nabla^2_{AY} \texttt{f}(\mb{A}, \mb{b}, \mb{y})
&=
\begin{bmatrix}
\begin{bmatrix}
2a_{00} y_0 + a_{01} y_1 + a_{02} y_2 & a_{00} y_1 & a_{00} y_2 \\
2a_{10} y_0 + a_{11} y_1 + a_{12} y_2 & a_{10} y_1 & a_{10} y_2
\end{bmatrix} \\ \\
%
\begin{bmatrix}
a_{01} y_0 & a_{00} y_0 + 2a_{01} y_1 + a_{02} x_2 & a_{01} y_2 \\
a_{11} y_0 & a_{10} y_0 + 2a_{11} y_1 + a_{12} x_2 &  a_{11} y_2
\end{bmatrix} \\ \\
%
\begin{bmatrix}
a_{02} y_0 & a_{02} y_1 & a_{00} y_0 + a_{01} y_1 + 2a_{02} y_2 \\
a_{12} y_0 & a_{12} y_1 & a_{10} y_0 + a_{11} y_1 + 2a_{12} y_2
\end{bmatrix}
\end{bmatrix}
-
\begin{bmatrix}
\begin{bmatrix}
b_0 & 0 & 0 \\
b_1 & 0 & 0
\end{bmatrix} \\ \\
\begin{bmatrix}
0 & b_0 & 0\\
0 & b_1 & 0
\end{bmatrix} \\ \\
\begin{bmatrix}
0 & 0 & b_0\\
0 & 0 & b_1
\end{bmatrix}
\end{bmatrix} \nonumber \\
&=
\begin{bmatrix}
\begin{bmatrix}
a_{00} y_0 & a_{00} y_1 & a_{00} y_2 \\
a_{10} y_0 & a_{10} y_1 & a_{10} y_2
\end{bmatrix}\\ \\
%
\begin{bmatrix}
a_{01} y_0 & a_{01} y_1 & a_{01} y_2 \\
a_{11} y_0 & a_{11} y_1 & a_{11} y_2
\end{bmatrix}\\ \\
%
\begin{bmatrix}
a_{02} y_0 & a_{02} y_1 & a_{02} y_2 \\
a_{12} y_0 & a_{12} y_1 & a_{12} y_2
\end{bmatrix}
\end{bmatrix}
+
\begin{bmatrix}
\begin{bmatrix}
a_{00} y_0 + a_{01} y_1 + a_{02} y_2 & 0 & 0 \\
a_{10} y_0 + a_{11} y_1 + a_{12} y_2 & 0 & 0
\end{bmatrix}\\ \\
%
\begin{bmatrix}
0 & a_{00} y_0 + a_{01} y_1 + a_{02} y_2 & 0 \\
0 & a_{10} y_0 + a_{11} y_1 + a_{12} y_2 & 0
\end{bmatrix}\\ \\
%
\begin{bmatrix}
0 & 0 & a_{00} y_0 + a_{01} y_1 + a_{02} y_2 \\
0 & 0 & a_{10} y_0 + a_{11} y_1 + a_{12} y_2
\end{bmatrix}
\end{bmatrix}
-
\begin{bmatrix}
\begin{bmatrix}
b_0 & 0 & 0 \\
b_1 & 0 & 0
\end{bmatrix} \\ \\
\begin{bmatrix}
0 & b_0 & 0\\
0 & b_1 & 0
\end{bmatrix} \\ \\
\begin{bmatrix}
0 & 0 & b_0\\
0 & 0 & b_1
\end{bmatrix}
\end{bmatrix} \nonumber \\
&=
\begin{bmatrix}
\begin{bmatrix}
a_{00} \\
a_{10}
\end{bmatrix} \mb{y}^T \\ \\
\begin{bmatrix}
a_{01} \\
a_{11}
\end{bmatrix} \mb{y}^T \\ \\
\begin{bmatrix}
a_{02} \\
a_{12}
\end{bmatrix} \mb{y}^T
\end{bmatrix}
+
\begin{bmatrix}
\mb{A} \mb{y} & \mb{0} & \mb{0} \\
\mb{0} & \mb{A} \mb{y} & \mb{0} \\
\mb{0} & \mb{0} & \mb{A} \mb{y}
\end{bmatrix}
-
\begin{bmatrix}
\mb{b} & \mb{0} & \mb{0} \\
\mb{0} & \mb{b} & \mb{0} \\
\mb{0} & \mb{0} & \mb{b}
\end{bmatrix} \nonumber \\
&=
\begin{bmatrix}
\mb{A}_0 \mb{y}^T \\
\mb{A}_1 \mb{y}^T \\
\mb{A}_2 \mb{y}^T
\end{bmatrix}
+
\begin{bmatrix}
\mb{A} \mb{y} - \mb{b} & \mb{0} & \mb{0} \\
\mb{0} & \mb{A} \mb{y} - \mb{b} & \mb{0} \\
\mb{0} & \mb{0} & \mb{A} \mb{y} - \mb{b}
\end{bmatrix}\ ,
\end{align}
where the left term can be written as $\mb{A}_i$ for all $i \in \mc{N} = \{1, ..., n\}$ and
the right term assigns $(\mb{A} \mb{y} - \mb{b})$ to the diagonal of $\nabla^2_{AY} \texttt{f}(\mb{A}, \mb{b}, \mb{y})$ which is $\mc{B}$ in Eq.~\eqref{eq:hessian_less}.
One can represent it as
\begin{align}
    \mc{B}_{ij} &= \mb{A}_i \mb{y}^T_j, \quad \forall i, j \in \mc{N}\ ,\\
    \mc{B}_{ii} &\leftarrow \mc{B}_{ii} + (\mb{A} \mb{y} - \mb{b}), \quad \forall i \in \mc{N}\ .
\end{align}
One can also obtain the same structure from different examples.
It is clear that forming $\mc{B}$ using $\mb{A}$, $\mb{b}$, and $\mb{y}$ in this way avoids calculating the zero vectors and matrices, and thus, becoming more efficient than using Jacobian to calculate every element of such a spare matrix.
To this end, we arrive at the result in Prop.~\ref{prop:less_hessian}. $\square$

\clearpage
Meanwhile, we show an implementation example comparing the automatic Jacobian (without exploited structures) and ours with the exploited structure in Listing~\ref{lst:second_order}.
%
\begin{lstlisting}[basicstyle=\footnotesize\ttfamily,linewidth=16cm,label={lst:second_order},caption={\textit{\normalsize \hspace{8mm}An implementation example of exploited Hessian structure for LESS.}\\}]
import torch
from torch.autograd import grad
from torch.autograd.functional import jacobian as J

def obj_fnc(A, u, b):
    Au_b = torch.einsum('bmn,bn->bm', A, u) - b
    loss = 0.5 * torch.einsum('bm,bm->b', Au_b, Au_b)

    return loss.sum()

# ==== Autogradient
def dfdy_auto_fnc(A, u, b):
    return grad(obj_fnc(A, u, b), (u), create_graph=True)[0]

def dffdAy_auto_fnc(A, u, b):
    m, n = A.shape[1:3]
    D = []

    with torch.enable_grad():
        for A_p, u_p, b_p in zip(A, u, b):
            A_p = A_p.view(1, m, n)
            u_p = u_p.view(1, n)
            b_p = b_p.view(1, m)
            D.append(J(lambda A: dfdy_auto_fnc(A, u_p, b_p), (A_p)))

        D = torch.cat(D, dim=0)

    return D

# ==== Structure exploited
def dfdy_fnc(A, u, b):
    Au_b = torch.einsum('bmn,bn->bm', A, u) - b

    return torch.einsum('bmn,bm->bn', A, Au_b)

def dffdAy_fnc(A, u, b):
    batch, m, n = A.shape
    D = A.new_zeros(batch, n, m, n)

    with torch.no_grad():
        for i in range(n):
            D_1 = torch.einsum('bm,bn->bmn', A[:, :, i], u)
            D_2 = torch.einsum('bmn,bn->bm', A, u) - b
            D[:, i] = D_1
            D[:, i, :, i] += D_2

    return D
\end{lstlisting}

% -------------------------------------------------------------------
\subsection{Proof for Proposition~\ref{prop:ied_gradient}}
For the principal matrix features, the implicit differentiation is on the eigenvector associated with the largest eigenvalue as
%
\begin{equation}
\nabla_X L(\mathbf{y})
= \nabla_Y L(\mathbf{y}) \nabla_X \mathbf{y}
= \nabla_Y L(\mathbf{y}) 
\left( \mc{H}\inv \mc{A}^{T} \left( \mc{A} \mc{H}\inv \mc{A}^T \right)\inv \mc{A} \mc{H}\inv - \mc{H}\inv \right) \mc{B}
= \mc{K} \mc{B} \in \mathbb{R}^{m \times m}\ ,
\end{equation}
%
where $\mc{K} \in \mathbb{R}^{m}$ is a vector given the upper-level loss $L \in \mathbb{R}$ from the bi-level problem.
Then,
%
\begin{equation}
\label{eq:prop_34_app}
\nabla_X L(\mathbf{y}) = -\mc{K} \mathbf{y}^T - \mathbf{y} \mc{K}^T\ .
\end{equation}

We find the matrix structure by taking the example of $m=2$, where $\mc{K}=[ k_0, k_1 ]^T$ and $\mb{y} = [ y_0, y_1 ]^T$.
Then,
%
\begin{equation}
\mc{B}
= \mc{B}_0 + \mc{B}_1
= \begin{bmatrix}
-y_0 & -y_1 & 0 & 0 \\
0 & 0 & -y_0 & -y_1
\end{bmatrix}
+
\begin{bmatrix}
-y_0 & 0 & -y_1 & 0 \\
0 & -y_0 & 0 & -y_1
\end{bmatrix}
\end{equation}
%
and
%
\begin{equation}
\mc{K} \mc{B}
= \mc{K} \mc{B}_0 + \mc{K} \mc{B}_1 =
\begin{bmatrix}
-k_0 y_0 & -k_0 y_1 \\
-k_1 y_0 & -k_1 y_1
\end{bmatrix}
+
\begin{bmatrix}
-k_0 y_0 & -k_1 y_0 \\
-k_0 y_1 & -k_1 y_1
\end{bmatrix}
=
-\begin{bmatrix}
k_0\\
k_1
\end{bmatrix}
%
\begin{bmatrix}
y_0 & y_1
\end{bmatrix}
-
\begin{bmatrix}
y_0\\
y_1
\end{bmatrix}
%
\begin{bmatrix}
k_0 & k_1
\end{bmatrix}
=
- \mc{K} \mb{y}^T - \mb{y} \mc{K}^T\ ,
\end{equation}
%
where $\mc{K} \mc{B}_0$ is the outer-product of $\mc{K}$ and $(-\mb{y})$, both $\mc{K} \mc{B}_1$ and $\mc{K} \mc{B}_0$ will be reshaped to $m \times m$ matrices.
One can also apply the same rule to different examples, which will have the same
structure in Eq.~\eqref{eq:prop_34_app}.
Clearly, using $\mb{y} \in \mathbb{R}^{m}$ for $\nabla_X L(\mathbf{y})$ is much more efficient and requires much less memory than using $\mc{B} \in \mathbb{R}^{m \times (m \times m)}$.
%
To this end, we prove Prop.~\ref{prop:ied_gradient}. $\square$

% -------------------------------------------------------------------
\subsection{More Results of LESS Optimization}
\begin{figure*}[!ht]
\centering
\includegraphics[width=\textwidth]{figures/case_study_more}
\vspace{-5mm}
\caption{\textit{Illustration of the iterative optimization in LESS.
Each row is a sample with 6 methods.
``{\color{blue}\textbf{blue}} dash": the least squares function;
``{\color{red}\textbf{red}} solid": the sphere constraints;
``{\color{green}\textbf{green}} solid": solution updates before and after $\texttt{Proj}_{\texttt{RM}}()$;
``{\color{green}\textbf{green}} dot-dash": the least squares function with the final solution;
``{\color{red}\textbf{red}} dot": the initial solution;
``{\color{black}\textbf{black}} star": the final solution, not always the optimal.
The least squares function with the exact solution ({\color{green}\textbf{green}} dot-dash) should be tangent to the sphere ({\color{red}\textbf{red}} solid).
Our ``PGD+RM+BLSOne" and ``PGD+RM+TWD" require much fewer iterations than the others for the optimal solution with comparable fixed point distance (FPD).}}
\label{fig:case_study_full}
\end{figure*}
%
In addition to Fig.~\ref{fig:case_study_short}, Fig.~\ref{fig:case_study_full} contains the full list of the compared methods.
The inputs for Fig.~\ref{fig:case_study_short} are
%
\begin{equation}
\mb{A} =
\begin{bmatrix}
0.569525 & -1.254572 \\
0.414020 & 0.124439
\end{bmatrix}\quad \text{and}
\quad
\mb{b} =
\begin{bmatrix}
-1.583332\\
-0.286124
\end{bmatrix}\ .
\end{equation}
%
In this case, the vanilla PGD requires more iterations than our ``PGD+RM+BLSOne" and ``PGD+RM+TWD".
Statistically, from Table~\ref{tb:table_1000}, when comparing ``PGD" with ``PGD+RM", all the metrics are better than ``PGD".
With ``+BLSOne" and ``+TWD", these can be further improved.

\begin{table*}[!ht]
\centering
\caption{\textit{Effectiveness evaluation on 1,000 random data sampled from $\mathcal{N}(0, 1)$ normal distribution.
\textbf{Bold}: the best, \underline{underline}: the second best.
``PGD": projected gradient descent (maximum 100 iterations),
``RM": Riemannian manifold,
``BLS": backtracking line search ($\alpha$=0.5 and $\beta$=0.8),
``TWD": tangent weight decay on Riemannian manifold ($\beta$=0.9),
``DW": direction weight,
``In" and ``Out": the number of failed cases inner and outer of the sphere respectively, where the case is regarded as failed when the solution update reaches the maximum 100 iterations,
``Imp.": the number of cases with energy no greater than SciPy.
Problem sizes, m$\times$n for A, are 2$\times$2 (282 inner and 718 outer), 64$\times$32 (555 inner and 445 outer), and 1024$\times$256 (all inner).}
}
\vspace{1mm}
\label{tb:table_1000}
\centering
\setlength{\tabcolsep}{3pt}
\resizebox{\textwidth}{!}{
\begin{tabular}{ccccccc|rrrr|rrrr|rrrr}
\hline
\multirow{2}{*}{SciPy} & \multirow{2}{*}{PGD} & \multirow{2}{*}{RM} & \multicolumn{2}{c}{BLS} & \multirow{2}{*}{TWD} & \multirow{2}{*}{DW} & \multicolumn{4}{c|}{2$\times$2} & \multicolumn{4}{c|}{64$\times$32} & \multicolumn{4}{c}{1024$\times$256} \\
& & & $\eta=1$ & $\eta=\nabla$ & & & In$\downarrow$ & Out$\downarrow$ & Imp.$\uparrow$ & MRE$\downarrow$ & In$\downarrow$ & Out$\downarrow$ & Imp.$\uparrow$ & MRE$\downarrow$ & In$\downarrow$ & Out$\downarrow$ & Imp.$\uparrow$ & MRE$\downarrow$ \\
\hline
\checkmark & & & & & & & \multicolumn{1}{c}{-} & \multicolumn{1}{c}{-} & \multicolumn{1}{c}{-} & \multicolumn{1}{c|}{-} & \multicolumn{1}{c}{-} & \multicolumn{1}{c}{-} & \multicolumn{1}{c}{-} & \multicolumn{1}{c|}{-} & \multicolumn{1}{c}{-} & \multicolumn{1}{c}{-} & \multicolumn{1}{c}{-} & \multicolumn{1}{c}{-} \\

& \checkmark & & & & & & \underline{172} & 353 & 389 & 10.84 & 548 & 440 & 52 & 0.15 & 1000 & \multicolumn{1}{c}{-} & 0 & 1.35 \\

& \checkmark & & & & & \checkmark & 179 & 438 & 37 & 4.89 & 474 & 363 & 63 & 0.07 & 1000 & \multicolumn{1}{c}{-} & 0 & 0.47 \\

& \checkmark & \checkmark & & & & & \textbf{0} & 118 & 844 & 4.30 & 30 & \textbf{0} & \underline{980} & \textbf{0.00} & \underline{4} & \multicolumn{1}{c}{-} & \textbf{1000} & \textbf{-0.03} \\

\rowcolor{lightgray} & \checkmark & \checkmark & \checkmark & & & & \textbf{0} & \underline{14} & \textbf{938} & \underline{0.14} & \textbf{6} & \textbf{0} & \textbf{994} & \textbf{0.00} & \underline{4} & \multicolumn{1}{c}{-} & \textbf{1000} & \textbf{-0.03} \\

& \checkmark & \checkmark & \checkmark & & & \checkmark & 175 & 421 & 40 & 3.51 & 285 & \underline{142} & 286 & \underline{0.01} & 984 & \multicolumn{1}{c}{-} & \underline{999} & \textbf{-0.03} \\

& \checkmark & \checkmark & & \checkmark & & & \textbf{0} & 108 & \underline{853} & 4.21 & \underline{7} & \textbf{0} & \textbf{994} & \textbf{0.00} & \textbf{2} & \multicolumn{1}{c}{-} & \textbf{1000} & \textbf{-0.03} \\

& \checkmark & \checkmark & & \checkmark & & \checkmark & 204 & 536 & 17 & 1.19 & 472 & 355 & 92 & 0.04 & 1000 & \multicolumn{1}{c}{-} & 0 & \underline{0.42} \\

\rowcolor{lightgray} & \checkmark & \checkmark & & & \checkmark & & \textbf{0} & \textbf{0} & 806 & \textbf{0.00} & 29 & \textbf{0} & 954 & \textbf{0.00} & \underline{4} & \multicolumn{1}{c}{-} & \textbf{1000} & \textbf{-0.03} \\

& \checkmark & \checkmark & & & \checkmark & \checkmark & 204 & 533 & 15 & 3.12 & 466 & 352 & 90 & 0.04 & 1000 & \multicolumn{1}{c}{-} & 0 & \underline{0.42} \\
\hline
\end{tabular}
}
\end{table*}

\clearpage
% -------------------------------------------------------------------
\subsection{Evaluation of IED Precision and Computational Requirements}

% ----
\subsubsection{Symmetric and Nonsymmetric Data Sampling}
\label{sec:sampling}
\textbf{Sampling Distributions}. Our evaluated data $\mb{A}$ are sampled from \textit{the standard Gaussian distribution} $\mc{N}(0, 1)$, \textit{uniform distribution} in $[0, 1)$, \textit{von Mises distribution} in $\mc{V}(0, 1)$ \cite{vonmise}, and \textit{random choices} from [0.0, 10.0), with absolute values, where $\mb{A} \leftarrow \mb{A}+\mb{A}^T$ is applied for symmetry.
We use $\texttt{numpy.random.randn}()$, $\texttt{numpy.random.uniform}()$, $\texttt{numpy.random.vonmises}(0, 1)$, and $\texttt{numpy.random.choice}(10.0)$, respectively, from the NumPy library \cite{harris2020array}.
They achieve similar numerical ranges as that from the Gaussian distribution in Figs.~\ref{fig:ied_symmetric_precision}-\ref{fig:ied_nonsymmetric_precision}.
Hyperparameters required by the sampling methods are not limited to the settings in this work, such as the mean and variance in the von Mises distribution and the value upper bound in the random choices.

% ----
\begin{figure}[!ht]
\centering
\begin{subfigure}{0.24\textwidth}
    \includegraphics[width=\textwidth]{figures/ied_data/solve_symmetric_stopFalse_float32_eigen_gap_gpu.png}
    \caption{ED, float32}
\end{subfigure}
\begin{subfigure}{0.24\textwidth}
    \includegraphics[width=\textwidth]{figures/ied_data/solve_symmetric_stopFalse_float64_eigen_gap_gpu.png}
    \caption{ED, float64}
\end{subfigure}
\begin{subfigure}{0.24\textwidth}
    \includegraphics[width=\textwidth]{figures/ied_data/solve_symmetric_stopFalse_float32_fp_gap_gpu.png}
    \caption{FPD, float32}
\end{subfigure}
\begin{subfigure}{0.24\textwidth}
    \includegraphics[width=\textwidth]{figures/ied_data/solve_symmetric_stopFalse_float64_fp_gap_gpu.png}
    \caption{FPD, float64}
\end{subfigure}
\vspace{-3mm}
\caption{\textit{Symmetric $\mb{A}$ sampled from Gaussian distribution with activation, precision evaluation with eigen distance and FPD.}}
\label{fig:ied_symmetric_precision}
\end{figure}

% ----
\begin{figure}[!ht]
\centering
\begin{subfigure}{0.24\textwidth}
    \includegraphics[width=\textwidth]{figures/ied_data/solve_nonsymmetric_stopFalse_float32_eigen_gap_gpu.png}
    \caption{ED, float32}
\end{subfigure}
\begin{subfigure}{0.24\textwidth}
    \includegraphics[width=\textwidth]{figures/ied_data/solve_nonsymmetric_stopFalse_float64_eigen_gap_gpu.png}
    \caption{ED, float64}
\end{subfigure}
\begin{subfigure}{0.24\textwidth}
    \includegraphics[width=\textwidth]{figures/ied_data/solve_nonsymmetric_stopFalse_float32_fp_gap_gpu.png}
    \caption{FPD, float32}
\end{subfigure}
\begin{subfigure}{0.24\textwidth}
    \includegraphics[width=\textwidth]{figures/ied_data/solve_nonsymmetric_stopFalse_float64_fp_gap_gpu.png}
    \caption{FPD, float64}
\end{subfigure}
\vspace{-3mm}
\caption{\textit{Nonsymmetric $\mb{A}$ sampled from Gaussian distribution with activation, precision evaluation with eigen distance and FPD.}}
\label{fig:ied_nonsymmetric_precision}
\end{figure}

% ----
\begin{figure}[!ht]
\centering
\begin{subfigure}{0.24\textwidth}
    \includegraphics[width=\textwidth]{figures/ied_data/resnet/pinverse_symmetric_stopFalse_float32_gaussian_resnet50_numseeds100_eigen_gap_gpu.png}
    \caption{Symmetric $\mb{A}$, ED}
\end{subfigure}
\begin{subfigure}{0.24\textwidth}
    \includegraphics[width=\textwidth]{figures/ied_data/resnet/pinverse_symmetric_stopFalse_float32_gaussian_resnet50_numseeds100_fp_gap_gpu.png}
    \caption{Symmetric $\mb{A}$, FPD}
\end{subfigure}
\begin{subfigure}{0.24\textwidth}
    \includegraphics[width=\textwidth]{figures/ied_data/resnet/pinverse_nonsymmetric_stopFalse_float32_gaussian_resnet50_numseeds100_eigen_gap_gpu.png}
    \caption{Nonsymmetric $\mb{A}$, ED}
\end{subfigure}
\begin{subfigure}{0.24\textwidth}
    \includegraphics[width=\textwidth]{figures/ied_data/resnet/pinverse_nonsymmetric_stopFalse_float32_gaussian_resnet50_numseeds100_fp_gap_gpu.png}
    \caption{Nonsymmetric $\mb{A}$, FPD}
\end{subfigure}
\vspace{-3mm}
\caption{\textit{Evaluation on 3 data sizes within the memory capacity as it requires a $(2048 \times 8^2) \times m^2$ dimension fully-connected layer in ResNet50, where $m=\{32, 64, 128\}$. All data is in float32.
Both metrics are the less the better.}}
\label{fig:ied_precision_resnet}
\end{figure}

The feasibility lies in 1) the simulation of normalized outputs (such as logits or softmax probability) of neural networks with an absolute operation,
2) the requirement of positive largest eigenvalue for IFT in Eq.~\eqref{eq:ift_pi}.
For nonsymmetric sampling, as we test 10,000 random samples (each with batch size 5) under this setting,
the power iteration and simultaneous iteration algorithms can always achieve almost 0 eigen distance and 0 fixed-point distance.
In contrast, PyTorch \texttt{eigh}()\footnote{PyTorch \texttt{eigh}() was designed for symmetric $\mb{A}$ and PyTorch \texttt{eig}() for complex and general data for nonsymmetric $\mb{A}$, but we are 1) not aim at complex data and 2) to avoid using \texttt{eigh}() for symmetric matrix but rather PI or SI by investigating these results.}, denoted as ``AutoDiff", has a much larger eigen distance, particularly in float32.

\textbf{Simulation on ResNet50}. We further generate random samples by using \textit{a modified ResNet50} with pretrained weights, removing the adaptive average pooling layer and adjusting the last fully-connected layer to obtain $m^2$ out channels (instead of 1,000) which will then be reshaped to a $m \times m$ matrix, followed by the absolute operation and symmetry for $\mc{A}$ if required.
It follows: random data with size ($5 \times 3 \times 256 \times 256$) that is (batch $\times$ channel $\times$ height $\times$ width) sampled from $\mc{N}(0, 1)$ $\rightarrow$ \texttt{ResNet}() $\rightarrow$ reshape to $m \times m$ $\rightarrow$ absolute operation $\rightarrow$ ($\mb{A} \leftarrow \mb{A} + \mb{A}^T$) if symmetric $\mb{A}$ is required.
In Figs.~\ref{fig:ied_precision_resnet}, ``PI" and ``SI" still achieve better eigen solutions than ``AutoDiff" on both symmetric and nonsymmetric $\mb{A}$.

% ----
\subsubsection{Additional Experiments on Computational Requirements}
We evaluate the solution precision with eigen distance and fixed-point distance on symmetric and nonsymmetric $\mb{A}$ in Fig.~\ref{fig:ied_symmetric_precision} and Fig.~\ref{fig:ied_nonsymmetric_precision} respectively.
Additionally, the computational requirements include running time and GPU memory in Tables~\ref{tb:ied_symmetric_time_cpu}-\ref{tb:ied_nonsymmetric_memory} for both symmetric and nonsymmetric $\mb{A}$.
Although ``AutoDiff" achieves faster speed in most cases, its precision is inferior to ``PI" and ``SI", particularly in float32.
For the evaluation completion, we provide all of these results.

Meanwhile, since a solver achieves the same solutions given the same inputs, it consumes the same (at least quite similar) computational resources in the backward pass.
We modulate each solver with different backward propagation methods into individual layers and provide all the forward and backward results.

In all these tables, implicit differentiation without exploited structure (``J") and with our exploited structure (``E") clearly distinguishes the benefits of exploiting the Jacobian and Hessian matrices in both the running time and memory requirements.
For those with out-of-memory issues or extensive running time, we mark the results with ``-" in the tables.

% ----
\begin{table*}[!ht]
\centering
\small
\caption{\textit{Implicit eigen decomposition (IED) layer, CPU time (s) for symmetric $\mb{A}$.
``AutoDiff": PyTorch \texttt{eigh}() function;
``DDN": deep declarative network;
``IFT": implicit function theorem;
``unroll": unrolling the forward iteration for gradients via PyTorch autodiff mechanism;
``J": autodiff Jacobian without exploited structure;
``E": our implicit differentiation with exploited structure.
Our \colorbox{gray!50}{best suggestions} are highlighted considering the overall solution precision and computational requirements.}
}
\label{tb:ied_symmetric_time_cpu}
\setlength{\tabcolsep}{11pt}
\resizebox{\textwidth}{!}{
\begin{tabular}{l|rr|rr|rr|rr}
\hline
\multicolumn{1}{c|}{\multirow{2}{*}{Method}} & \multicolumn{2}{c|}{5$\times$32} & \multicolumn{2}{c|}{5$\times$64} & \multicolumn{2}{c|}{5$\times$128} & \multicolumn{2}{c}{5$\times$256} \\ 
 & Forward & Backward & Forward & Backward & Forward & Backward & Forward & Backward
\csvreader[head to column names]{figures/ied_data/solve_symmetric_stopFalse_float32/time_cpu.csv}{}
{\\\hline\csvcoli & \csvcolii & \csvcoliii & \csvcoliv & \csvcolv & \csvcolvi & \csvcolvii & \csvcolviii & \csvcolix}
\\\hline
\end{tabular}
}
\end{table*}

% ----
\begin{table*}[!ht]
\centering
\small
\caption{\textit{Implicit eigen decomposition (IED) layer, GPU time (s) for symmetric $\mb{A}$.
``AutoDiff": PyTorch \texttt{eigh}() function;
``DDN": deep declarative network;
``IFT": implicit function theorem;
``unroll": unrolling the forward iteration for gradients via PyTorch autodiff mechanism;
``J": autodiff Jacobian without exploited structure;
``E": our implicit differentiation with exploited structure.
Our \colorbox{gray!50}{best suggestions} are highlighted considering the overall solution precision and computational requirements.}
}
\label{tb:ied_symmetric_time_gpu}
\setlength{\tabcolsep}{2pt}
\resizebox{\textwidth}{!}{
\begin{tabular}{l|rr|rr|rr|rr|rr|rr}
\hline
\multicolumn{1}{c|}{\multirow{2}{*}{Method}} & \multicolumn{2}{c|}{5$\times$32} & \multicolumn{2}{c|}{5$\times$64} & \multicolumn{2}{c|}{5$\times$128} & \multicolumn{2}{c|}{5$\times$256} & \multicolumn{2}{c|}{5$\times$512} & \multicolumn{2}{c}{5$\times$1024} \\ 
 & Forward & Backward & Forward & Backward & Forward & Backward & Forward & Backward & Forward & Backward & Forward & Backward
\csvreader[head to column names]{figures/ied_data/solve_symmetric_stopFalse_float32/time_gpu.csv}{}
{\\\hline\csvcoli & \csvcolii & \csvcoliii & \csvcoliv & \csvcolv & \csvcolvi & \csvcolvii & \csvcolviii & \csvcolix & \csvcolx & \csvcolxi & \csvcolxii & \csvcolxiii}
\\\hline
\end{tabular}
}
\end{table*}

% ----
\begin{table*}[!ht]
\centering
\small
\caption{\textit{Implicit eigen decomposition (IED) layer, GPU memory (MB) for symmetric $\mb{A}$.
``AutoDiff": PyTorch \texttt{eigh}() function;
``DDN": deep declarative network;
``IFT": implicit function theorem;
``unroll": unrolling the forward iteration for gradients via PyTorch autodiff mechanism;
``J": autodiff Jacobian without exploited structure;
``E": our implicit differentiation with exploited structure.
Our \colorbox{gray!50}{best suggestions} are highlighted considering the overall solution precision and computational requirements.}
}
\label{tb:ied_symmetric_memory}
\setlength{\tabcolsep}{0.7pt}
\resizebox{\textwidth}{!}{
\begin{tabular}{l|rr|rr|rr|rr|rr|rr}
\hline
\multicolumn{1}{c|}{\multirow{2}{*}{Method}} & \multicolumn{2}{c|}{5$\times$32} & \multicolumn{2}{c|}{5$\times$64} & \multicolumn{2}{c|}{5$\times$128} & \multicolumn{2}{c|}{5$\times$256} & \multicolumn{2}{c|}{5$\times$512} & \multicolumn{2}{c}{5$\times$1024} \\ 
 & Forward & Backward & Forward & Backward & Forward & Backward & Forward & Backward & Forward & Backward & Forward & Backward
\csvreader[head to column names]{figures/ied_data/solve_symmetric_stopFalse_float32/memory_gpu.csv}{}
{\\\hline\csvcoli & \csvcolii & \csvcoliii & \csvcoliv & \csvcolv & \csvcolvi & \csvcolvii & \csvcolviii & \csvcolix & \csvcolx & \csvcolxi & \csvcolxii & \csvcolxiii}
\\\hline
\end{tabular}
}
\end{table*}

% ----
\begin{table*}[!ht]
\centering
\small
\caption{\textit{Implicit eigen decomposition (IED) layer, CPU time (s) for nonsymmetric $\mb{A}$.
``AutoDiff": PyTorch \texttt{eigh}() function;
``DDN": deep declarative network;
``IFT": implicit function theorem;
``unroll": unrolling the forward iteration for gradients via PyTorch autodiff mechanism;
``J": autodiff Jacobian without exploited structure;
``E": our implicit differentiation with exploited structure.
Our \colorbox{gray!50}{best suggestions} are highlighted considering the overall solution precision and computational requirements.}
}
\label{tb:ied_nonsymmetric_time_cpu}
\setlength{\tabcolsep}{11pt}
\resizebox{\textwidth}{!}{
\begin{tabular}{l|rr|rr|rr|rr}
\hline
\multicolumn{1}{c|}{\multirow{2}{*}{Method}} & \multicolumn{2}{c|}{5$\times$32} & \multicolumn{2}{c|}{5$\times$64} & \multicolumn{2}{c|}{5$\times$128} & \multicolumn{2}{c}{5$\times$256} \\ 
 & Forward & Backward & Forward & Backward & Forward & Backward & Forward & Backward
\csvreader[head to column names]{figures/ied_data/solve_nonsymmetric_stopFalse_float32/time_cpu.csv}{}
{\\\hline\csvcoli & \csvcolii & \csvcoliii & \csvcoliv & \csvcolv & \csvcolvi & \csvcolvii & \csvcolviii & \csvcolix }
\\\hline
\end{tabular}
}
\end{table*}

% ----
\begin{table*}[!ht]
\centering
\small
\caption{\textit{Implicit eigen decomposition (IED) layer, GPU time (s) for nonsymmetric $\mb{A}$.
``AutoDiff": PyTorch \texttt{eigh}() function;
``DDN": deep declarative network;
``IFT": implicit function theorem;
``unroll": unrolling the forward iteration for gradients via PyTorch autodiff mechanism;
``J": autodiff Jacobian without exploited structure;
``E": our implicit differentiation with exploited structure.
Our \colorbox{gray!50}{best suggestions} are highlighted considering the overall solution precision and computational requirements.}
}
\label{tb:ied_nonsymmetric_time_gpu}
\setlength{\tabcolsep}{2pt}
\resizebox{\textwidth}{!}{
\begin{tabular}{l|rr|rr|rr|rr|rr|rr}
\hline
\multicolumn{1}{c|}{\multirow{2}{*}{Method}} & \multicolumn{2}{c|}{5$\times$32} & \multicolumn{2}{c|}{5$\times$64} & \multicolumn{2}{c|}{5$\times$128} & \multicolumn{2}{c|}{5$\times$256} & \multicolumn{2}{c|}{5$\times$512} & \multicolumn{2}{c}{5$\times$1024} \\ 
 & Forward & Backward & Forward & Backward & Forward & Backward & Forward & Backward & Forward & Backward & Forward & Backward
\csvreader[head to column names]{figures/ied_data/solve_nonsymmetric_stopFalse_float32/time_gpu.csv}{}
{\\\hline\csvcoli & \csvcolii & \csvcoliii & \csvcoliv & \csvcolv & \csvcolvi & \csvcolvii & \csvcolviii & \csvcolix & \csvcolx & \csvcolxi & \csvcolxii & \csvcolxiii}
\\\hline
\end{tabular}
}
\end{table*}

% ----
\begin{table*}[!ht]
\centering
\small
\caption{\textit{Implicit eigen decomposition (IED) layer, GPU memory (MB) for nonsymmetric $\mb{A}$.
``AutoDiff": PyTorch \texttt{eigh}() function;
``DDN": deep declarative network;
``IFT": implicit function theorem;
``unroll": unrolling the forward iteration for gradients via PyTorch autodiff mechanism;
``J": autodiff Jacobian without exploited structure;
``E": our implicit differentiation with exploited structure.
Our \colorbox{gray!50}{best suggestions} are highlighted considering the overall solution precision and computational requirements.}
}
\label{tb:ied_nonsymmetric_memory}
\setlength{\tabcolsep}{0.7pt}
\resizebox{\textwidth}{!}{
\begin{tabular}{l|rr|rr|rr|rr|rr|rr}
\hline
\multicolumn{1}{c|}{\multirow{2}{*}{Method}} & \multicolumn{2}{c|}{5$\times$32} & \multicolumn{2}{c|}{5$\times$64} & \multicolumn{2}{c|}{5$\times$128} & \multicolumn{2}{c|}{5$\times$256} & \multicolumn{2}{c|}{5$\times$512} & \multicolumn{2}{c}{5$\times$1024} \\ 
 & Forward & Backward & Forward & Backward & Forward & Backward & Forward & Backward & Forward & Backward & Forward & Backward
\csvreader[head to column names]{figures/ied_data/solve_nonsymmetric_stopFalse_float32/memory_gpu.csv}{}
{\\\hline\csvcoli & \csvcolii & \csvcoliii & \csvcoliv & \csvcolv & \csvcolvi & \csvcolvii & \csvcolviii & \csvcolix & \csvcolx & \csvcolxi & \csvcolxii & \csvcolxiii}
\\\hline
\end{tabular}
}
\end{table*}